% Define a counter named 'exercisecounter' to track the exercise.
% The counter will be reset to zero each time the 'chapter' counter
% is incremented (when a new chapter is created for example).
% See pandoc_filter.py3
% The same for 'projectcounter' to track projects.
\newcounter{exercisecounter}[chapter]
\newcounter{projectcounter}[chapter]

\usepackage{graphicx}
\usepackage{wrapfig}

% Map Unicode emojis to PDF images. Based on coloremoji.sty
% See https://github.com/alecjacobson/coloremoji.sty
\usepackage{newunicodechar}
\newunicodechar{😈}{\text{\raisebox{-0.2em}{\includegraphics[height=1em]{img/emoji/evil.pdf}}}}
\newunicodechar{💣}{\text{\raisebox{-0.2em}{\includegraphics[height=1em]{img/emoji/bomb.pdf}}}}
\newunicodechar{⚠}{\text{\raisebox{-0.2em}{\includegraphics[height=1em]{img/emoji/alert.pdf}}}}


% Disabled because it enters in conflict with tikz
% The error says:
%
%   ! Argument of \language@active@arg> has an extra }.
%   <inserted text>
%                   \par
%   l.86 ...in{tikzpicture}[>=latex',line join=bevel,]
%\usepackage[spanish]{babel}

\usepackage{draftwatermark}
\SetWatermarkColor[gray]{0.96}

% call \chapterauthor{aaa} *after* each chapter begin to name the author
% of it
\makeatletter
\newcommand{\chapterauthor}[1]{%
  {\parindent0pt\vspace*{-25pt}%
  \linespread{1.1}\large\scshape#1%
  \par\nobreak\vspace*{35pt}}
  \@afterheading%
}
\makeatother

\title{Guía de Taller}
\author{Martín Di Paola}

\hypersetup{
            pdftitle={Guía de Taller},
            pdfauthor={Martín Di Paola},
            pdfborder={0 0 0},
            breaklinks=true}

\newfontfamily\monacofont[Path=./fonts/]{monaco.ttf}

% Define some colors
\usepackage{xcolor}
\definecolor{darkgreen}{rgb}{0,.5,0}
\definecolor{darkblue}{rgb}{0,0,.7}

% Define some styles for the highlighting with 'listings'
\lstdefinestyle{normal}{                    %
   numbers=none,                            % enumerar las lineas
   showstringspaces=false,                  % NOTE: without this xelatex/xdvipdfmx breaks! (i don't know why)
   keywordstyle=\color{darkblue}\textbf,    % color de las keywords
   stringstyle=\color{red},                 % color de los strings
   commentstyle=\color{darkgreen},          % color de los comentarios
   basicstyle=\color{black}\monacofont,     % color del texto en general
   morecomment=[l][\color{magenta}]{\#},    % coloreamos las intrucciones del precompilador (todo lo que empieza con #)
   ndkeywords={NULL,nullptr,noexcept},      % definimos una nuevas keywords como NULL y nullptr
   ndkeywordstyle=\color{violet},           % y las nuevas keywords tendran este color
   frame=none,                              % simple, sin ningun marco o frame alrededor del codigo
   basewidth={0.55em,0.55em}                % tamano de las letras/lineas. usado para reducir el whitespace entre estas
}

\lstset{style=normal}

% Graph and tree diagrams
\usepackage{tikz}
\usetikzlibrary{snakes,arrows,shapes}
\usepackage{amsmath}
